\documentclass[10pt,a4paper]{amsart}
\usepackage[margin=19mm]{geometry}
\usepackage{amsmath,amssymb,mathtools}
\usepackage[scr=rsfs,cal=euler]{mathalfa}
\usepackage{xfrac}
\usepackage[unicode,hidelinks,bookmarks=false]{hyperref}
\newcommand{\ie}{i.e.\ }
\newcommand{\cf}{cf.\ }
\newcommand{\resp}{respectively\ }
\newcommand{\Subsection}{\subsection}
\providecommand{\nobreakdash}{\nobreak\hbox{-}}
\setlength{\emergencystretch}{2em}
\multlinegap0pt
\usepackage{mathtools}
\usepackage{tikz}
\usepackage{xcolor}
\usepackage{amssymb}
\usepackage{float}
\usepackage{algorithm}\usepackage{algpseudocode}
\newtheorem{lemma}{Lemma}
\newtheorem{theorem}{Theorem}
\newcommand{\boff}{b_h^{\off}} % bonus  offline
\newcommand{\bont}{b_h^{t,\on}} % bonus t online
\newcommand{\highv}{U} % theoretical upper bound on V
\newcommand{\lowv}{W} % theoretical lower bound on V
\newcommand{\off}{\mathrm{off}} % offline
\newcommand{\ohighv}{\overline{V}} % offline V upper bound
\newcommand{\olowv}{\underline{V}} % offline V lower bound
\newcommand{\on}{\mathrm{on}} % online
\title{Learning Upper Lower Value Envelopes to Shape Online RL: A Principled Approach}
\author{Sebastian Reboul, H\'el\`ene Halconruy}
\date{Source: arXiv:2510.19528v2 (CC BY 4.0)}
\begin{document}
\maketitle

\noindent\emph{Source and licence.} Learning Upper Lower Value Envelopes to Shape Online RL: A Principled Approach. Version arXiv:2510.19528v2 (CC BY 4.0), distributed by arXiv under the Creative Commons Attribution 4.0 International licence, http://creativecommons.org/licenses/by/4.0/, as recorded in the arXiv licence metadata for this exact version and shown on the abstract page https://arxiv.org/abs/2510.19528v2. Full licence notice: \texttt{licenses/arxiv-2510.19528-CC-BY-4.0.txt}.\par\smallskip

\noindent\textit{Editorial scope.} Counted unit: the article's lemma lem:on\_bonus\_validity (source0.tex lines 1657--1669). The article's complete original proof of it is delivered with it (source0.tex lines 1671--1767, one \texttt{proof} environment of 97 lines). No mathematics is written, repaired or re-derived: the statement and the proof are the article's own lines, in the article's own notation, and no retained line is substituted or re-flowed.  Retained uncounted as the source-local context the proof needs: lines 1121--1127; lines 1364--1422.  Nothing was omitted from any retained span, so the package carries no omission record.  No retained line contains a citation, so the wrapper carries no bibliography and no bibliography entry.  No retained line contains a figure environment, an image-inclusion command or drawing code, so no image is delivered and the package declares no asset dependency.  Editorial formatting only, in addition: an AMS \texttt{amsart} preamble that restates the article's own declarations and macros used by the retained lines, namely the environments \texttt{lemma}, \texttt{theorem} on separate counters, together with the standard packages those lines need.  The retained environments print in this document as Lemma 1, Theorem 1 in order of appearance, whereas the article prints the same items with the numbering of its own full document.  The complete native source of the article as distributed in the arXiv e-print for this exact version is bundled unchanged as \texttt{sources/187579/source0.tex}.
\begin{lemma}[Offline Bonus Validity] 
\label{lem:off_bonus_validity} 
With probability $\ge 1 - \delta$, for every $h, s, a$, the bonuses provide valid confidence intervals: 
\begin{equation}\label{eq:off_bonus_validity} 
|(P^{\star}-\widehat{P}^{(h)})\olowv_{h+1} |\le \boff(s,a) \quad\text{and}\quad |(P^{\star}-\widehat{P}^{(h)})\ohighv_{h+1}| \le \boff(s,a). 
\end{equation} 
\end{lemma}

\begin{theorem}[Conditionally Independent Empirical Bernstein]
\label{thm:cond_empirical_bernstein_final}
Let $(\Omega,\mathcal{F},\mathbb{P})$ be a probability space and let $\mathcal{G}\subseteq\mathcal{F}$ be a sub-$\sigma$-algebra.
Let $(\mathcal{X},\mathcal{B})$ be a measurable space and let $X_1,\dots,X_n$ ($n\ge 2$) be i.i.d.\ $\mathcal{X}$-valued random variables such that $\sigma(X_1,\dots,X_n)$ is independent of $\mathcal{G}$.

Let $g:\Omega\times\mathcal{X}\to\mathbb{R}$ be $(\mathcal{G}\otimes\mathcal{B})$-measurable.
Let $\underline{g},\overline{g}:\Omega\to\mathbb{R}$ be $\mathcal{G}$-measurable, assume $\underline{g}\le \overline{g}$ almost surely, and define the (random) range length
\[
R:=\overline{g}-\underline{g}.
\]
Let $E\in\mathcal{G}$ be an event such that on $E$,
\[
\underline{g}(\omega)\le g(\omega,x)\le \overline{g}(\omega)
\qquad
\text{for all }x\in\mathcal{X}.
\]

For each $i=1,\dots,n$, define the random variable $Z_i:\Omega\to\mathbb{R}$ by
\[
Z_i(\omega):=g(\omega,X_i(\omega)).
\]
Let
\[
\overline{Z}(\omega):=\frac{1}{n}\sum_{i=1}^n Z_i(\omega),
\qquad
\widehat{\mathrm{Var}}(\mathbf Z)(\omega)
:=\frac{1}{n}\sum_{i=1}^n (Z_i(\omega)-\overline{Z}(\omega))^2,
\qquad
\mu_g(\omega):=\mathbb E[Z_1\mid \mathcal{G}](\omega).
\]
Then for any $\delta\in(0,1)$,
\[
\mathbb{P}\Bigl(
E\cap\Bigl\{
\bigl|\mu_g-\overline{Z}\bigr|
\le
2\sqrt{\frac{\widehat{\mathrm{Var}}(\mathbf Z)\ln(4/\delta)}{n}}
+
\frac{14\,R\,\ln(4/\delta)}{3n}
\Bigr\}
\Bigr)
\ge
(1-\delta)\,\mathbb{P}(E),
\]
equivalently,
\[
\mathbb{P}\Bigl(
\bigl\{
\bigl|\mu_g-\overline{Z}\bigr|
>
2\sqrt{\frac{\widehat{\mathrm{Var}}(\mathbf Z)\ln(4/\delta)}{n}}
+
\frac{14\,R\,\ln(4/\delta)}{3n}
\bigr\}\cap E
\Bigr)
\le
\delta\,\mathbb{P}(E).
\]
\end{theorem}

\begin{lemma}[Online Bonus Validity]
\label{lem:on_bonus_validity}
Define :
\begin{equation}\label{eq:def_Eon}
\mathcal{E}^{\mathrm{On}}_\delta=\{(t,s,a,h) \in [T]\times\mathcal{S} \times \mathcal{A} \times [H]:\eqref{eq:on_bonus_validity}\,\mathrm{holds}\}.
\end{equation}
where for every $t, h, s, a$, \eqref{eq:on_bonus_validity} is
\begin{equation}\label{eq:on_bonus_validity}
| \left \langle P^{\star}_{h,s,a}-\widehat{P}_{h,s,a}^{t},V_{h+1}^{\star} \right \rangle| \le \bont(s,a) \quad \text{and} \quad | \left \langle P_{h,s,a}^{\star}-\widehat{P}_{h,s,a}^{t}, \highv_{h+1} \right \rangle| \le \bont(s,a).
\end{equation}
Then we have that
\[ \mathbb{P}(\mathcal{E}^{\mathrm{Off}}_\delta \cap \mathcal{E}^{\mathrm{On}}_\delta) \ge 1 - 2\delta.\]
\end{lemma}

\begin{proof}
Define $(X_{k,h}^{s,a})_{(s,a)\in\mathcal S\times\mathcal A,\ h\in[H],\ k\in[T]}$ taking values in $\mathcal S$ and independent across $s,a,h,k$, such as
  \[
  \mathbb{P}(X_{k,h}^{s,a}=s')=P_h(s'\mid s,a)\ \text{ for all }s'\in\mathcal S.
  \]
Our assumption on the model is that the offline dataset is independent of the array $(X_{k,h}^{s,a})$, and that the online dynamics are \emph{coupled} to the pre-generated transitions by the following protocol:
\[
\textit{Whenever $(s,a)$ is visited at step $h$ for the $k$-th time during the online phase, then the next state is } X_{k,h}^{s,a}.
\]
For $(s,a,h)$ and $t\in[T]$, let
\[
\mathcal T_h^t(s,a)=\{t'\le t:\ S_h^{t'}=s,\ A_h^{t'}=a\},\qquad
N_h^t(s,a)=|\mathcal T_h^t(s,a)|.
\]
Order $\mathcal T_h^t(s,a)$ as $t'_1<\dots<t'_{N_h^t(s,a)}$ and write $S_i'=S_{h+1}^{t'_i}$. By the coupling,
\[
S_i'=X_{i,h}^{s,a}\qquad\text{for }i=1,\dots,N_h^t(s,a).
\]
The empirical kernel based on the first $n$ i.i.d. transitions at step $h$ is
\[
\widehat P_h^{[n]}(s'\mid s,a)=\frac1n\sum_{k=1}^n\mathbf 1_{\{X_{k,h}^{s,a}=s'\}},
\]
and the online empirical kernel is
\[
\widehat P_h^{t}(s'\mid s,a)=\frac{1}{N_h^t(s,a)}\sum_{t'\in\mathcal T_h^t(s,a)}\mathbf 1_{\{S_{h+1}^{t'}=s'\}}.
\]
By the coupling, for $n=N_h^t(s,a)\ge1$,
\[
\widehat P_h^{t}(\cdot\mid s,a)=\widehat P_h^{[n]}(\cdot\mid s,a).
\]

Fix $(h,s,a)$ and $n \in [T]$ with $n \ge 2$. We will apply  Theorem~\ref{thm:cond_empirical_bernstein_final} to two distinct choices for the target function $g: \Omega \times \mathcal{S} \to \mathbb{R}$:
\begin{enumerate}
    \item $g(\omega, s') = V_{h+1}^\star(s')$ (which is deterministic and thus $\mathcal{G}_K$-measurable).
    \item $g(\omega, s') = \highv_{h+1}(\omega, s')$ (which is a random variable computed from the offline dataset).
\end{enumerate}
We use $\mathcal{G} = \mathcal{G}_K$ and the event $E = \mathcal{E}^{\mathrm{Off}}_\delta$.
We use the coupled next-state variables $X_{1,h}^{s,a}, \dots, X_{n,h}^{s,a}$ as the sample $X_1, \dots, X_n$.
We define $\underline{g} \coloneqq \min_{z} \lowv_{h+1}(z)$ and $\overline{g} \coloneqq \max_{z} \highv_{h+1}(z)$, yielding range $R_{h+1} = \overline{g} - \underline{g}$.
We define $Z_i \coloneqq g(\omega, X_{i,h}^{s,a}())$. This establishes the direct parallel to the theorem's result terms:
    \begin{align*}
        \overline{Z}(\omega) &= \frac{1}{n}\sum_{i=1}^n g(\omega, X_{i,h}^{s,a}(\omega)) = \langle \widehat{P}_h^{[n]}(\cdot\mid s,a)(\omega), g(\omega,\cdot) \rangle, \\
        \mu_g(\omega) &= \langle P_{h,s,a}^\star(\cdot), g(\omega,\cdot) \rangle, \\
        \widehat{\mathrm{Var}}(\mathbf{Z}) &= \frac{1}{n}\sum_{i=1}^n (Z_i - \overline{Z})^2 = \widehat{\mathrm{Var}}_{\widehat{P}_h^{[n]}(\cdot\mid s,a)}[g]
    \end{align*}

Applying Theorem~\ref{thm:cond_empirical_bernstein_final} with confidence level $\delta'$, we substitute these equivalences directly into the concentration bound. We obtain that for each choice of $g$:
\[
\mathbb{P}\left(
\mathcal E_\delta^{\mathrm{Off}} \cap \left\{
\left| \langle P_{h,s,a}^\star-\widehat P_h^{[n]}(\cdot\mid s,a), g \rangle \right|
\le
2\sqrt{\frac{\widehat{\mathrm{Var}}_{\widehat P_h^{[n]}(\cdot\mid s,a)}[g]\ln(4/\delta')}{n}}
+
\frac{14 R_{h+1}\ln(4/\delta')}{3n}
\right\}
\right)
\ge
(1-\delta')\mathbb{P}(\mathcal E_\delta^{\mathrm{Off}}).
\]
Furthermore, on the event $\mathcal E_\delta^{\mathrm{Off}}$, the values of $g$ are confined to an interval of length $R_{h+1}$, which implies the deterministic bound $| \langle P_{h,s,a}^\star-\widehat P_h^{[n]}, g \rangle | \le R_{h+1}$.
 

The next step is to bound the biased variance $\widehat{\mathrm{Var}}_{\widehat{P}^{t}_h(\cdot|s,a)}[g]$.
Recall 
\[ M_h = \frac{1}{2}(\highv_{h+1} + \lowv_{h+1}), \quad D_h = \highv_{h+1} - \lowv_{h+1}, \]
the fact that $\lowv_{h+1} \le g \le \highv_{h+1}$ implies that $|g - M_{h+1}| \le \dfrac{D_{h+1}}{2}$. Therefore, decomposing 
\[ g = M_{h+1} + (g - M_{h+1}) \]
and applying the Minkowski inequality to variances we get
\begin{align*}
\sqrt{\widehat{\mathrm{Var}}_{\widehat{P}^{t}_h(\cdot|s,a)}[g]} 
    & \leq 
\sqrt{\widehat{\mathrm{Var}}_{\widehat{P}^{t}_h(\cdot|s,a)}[M_{h+1}]} + \sqrt{\widehat{\mathrm{Var}}_{\widehat{P}^{t}_h(\cdot|s,a)}[g - M_{h+1}]} \\
    & \leq \sqrt{\widehat{\mathrm{Var}}_{\widehat{P}^{t}_h(\cdot|s,a)}[M_{h+1}]} + \sqrt{\widehat{\mathbb{E}}_{\widehat{P}^{t}_h(\cdot|s,a)}[(g - M_{h+1})^2]} \\
    & \leq 
\sqrt{\widehat{\mathrm{Var}}_{\widehat{P}^{t}_h(\cdot|s,a)}[M_{h+1}]} + \frac{1}{2}\sqrt{\widehat{\mathbb{E}}_{\widehat{P}^{t}_h(\cdot|s,a)}[D_{h+1}^2]} \\
    & = \sigma_{h+1}^t(s,a).
\end{align*}

Combining with the previous step and taking the minimum with $R_{h+1}$ yields exactly the online bonus form $\bont(s,a)$ used in \eqref{eq:on_bonus_validity}.

Choose $\delta' = \delta/(2|\mathcal S||\mathcal A|HT)$ and take a union bound over all $h \in [H]$, $(s,a)\in\mathcal S\times\mathcal A$, $n \in [T]$, and both choices $g \in \{V_{h+1}^\star,\highv_{h+1}\}$.
Since Theorem \ref{thm:cond_empirical_bernstein_final} controls each deviation event on $\mathcal E_\delta^{\mathrm{Off}}$ with conditional failure probability at most $\delta'$, we obtain
\[
\mathbb{P}\bigl(\mathcal E_\delta^{\mathrm{Off}} \cap \mathcal E_\delta^{\mathrm{On}}\bigr)
\ge
(1-\delta)\mathbb{P}(\mathcal E_\delta^{\mathrm{Off}}).
\]
Finally, Lemma \ref{lem:off_bonus_validity} gives $\mathbb{P}(\mathcal E_\delta^{\mathrm{Off}})\ge 1-\delta$, hence
\[
\mathbb{P}\bigl(\mathcal E_\delta^{\mathrm{Off}} \cap \mathcal E_\delta^{\mathrm{On}}\bigr)
\ge
(1-\delta)^2
\ge
1-2\delta.
\]
\end{proof}

\end{document}
